\section{Problem Set 8: Sequences and Functions}

\begin{enumerate}
\item \textbf{Problem 1}\par
Decide which of the following assignments define functions:
\begin{enumerate}
\item each real number \(x\) is assigned \(x^2\);
\item each positive number \(x\) is assigned both \(\sqrt{x}\) and \(-\sqrt{x}\);
\item each student in a class is assigned the student's date of birth.
\end{enumerate}
Explain each decision.

\item \textbf{Problem 2}\par
For
\[
f:\mathbb{R}\to\mathbb{R},
\qquad
f(x)=x^2-4,
\]
state the domain and codomain, and determine the range.

\item \textbf{Problem 3}\par
Let
\[
f(x)=2x^2-3x+1.
\]
Compute \(f(0)\), \(f(-2)\), \(f(a)\), and \(f(x+h)\).

\item \textbf{Problem 4}\par
If
\[
g(x)=\frac{1}{x-3},
\]
determine the natural real domain of \(g\) and compute \(g(1)\).

\item \textbf{Problem 5}\par
The graph of a function contains the points
\[
(-2,3),\quad (0,-1),\quad (4,5).
\]
State the corresponding function values.

\item \textbf{Problem 6}\par
For
\[
f(x)=x^2-5x+6,
\]
find the \(y\)-intercept and all \(x\)-intercepts.

\item \textbf{Problem 7}\par
Determine the natural real domain of each function:
\[
f(x)=\sqrt{x+2},
\qquad
g(x)=\frac{x+1}{x^2-9},
\qquad
h(x)=\ln(4-x).
\]

\item \textbf{Problem 8}\par
Find the range of each function:
\[
f(x)=x^2+2,
\qquad
g(x)=-3x+1,
\qquad
h(x)=\sqrt{x}.
\]

\item \textbf{Problem 9}\par
Classify each function as constant, affine, quadratic, polynomial, rational, exponential, logarithmic, or trigonometric:
\[
3,\quad 2x-5,\quad x^2+1,\quad \frac{x+1}{x-2},\quad 4^x,\quad \ln x,\quad \sin x.
\]

\item \textbf{Problem 10}\par
Determine which of the functions
\[
f(x)=3x,\qquad g(x)=3x+2,\qquad h(x)=-x
\]
are linear in the sense \(f(x+y)=f(x)+f(y)\). Explain why an affine function need not be linear.

\item \textbf{Problem 11}\par
For
\[
f(x)=2^x
\quad\text{and}\quad
g(x)=\log_2 x,
\]
compute \(f(-2)\), \(f(3)\), \(g(1)\), \(g(8)\), and explain the relationship between the two functions.

\item \textbf{Problem 12}\par
Let
\[
a_n=3n-1,\qquad n\ge1.
\]
Write the first five terms and compute \(a_{20}\).

\item \textbf{Problem 13}\par
A sequence is defined by
\[
a_1=4,\qquad a_{n+1}=a_n-3.
\]
Write the first six terms.

\item \textbf{Problem 14}\par
Find an explicit formula for the arithmetic sequence
\[
5,\ 9,\ 13,\ 17,\ldots
\]

\item \textbf{Problem 15}\par
Find an explicit formula for the geometric sequence
\[
3,\ 6,\ 12,\ 24,\ldots
\]

\item \textbf{Problem 16}\par
Define recursively the sequence
\[
2,\ 6,\ 18,\ 54,\ldots
\]
and state the starting index.

\item \textbf{Problem 17}\par
For
\[
a_n=\frac{1}{n},
\]
list the first five graph points \((n,a_n)\). Explain why the graph should normally be drawn as separated points.

\item \textbf{Problem 18}\par
Compare the domains of
\[
a_n=n^2,\qquad n\in\mathbb{N},
\]
and
\[
f(x)=x^2,\qquad x\in\mathbb{R}.
\]
Which expressions \(a_{3/2}\) and \(f(3/2)\) are meaningful?

\item \textbf{Problem 19}\par
A graph rises on \((-3,1)\), falls on \((1,4)\), crosses the \(x\)-axis at \(-2\) and \(3\), and has a local maximum at \((1,5)\). Translate this information into statements about the function.

\item \textbf{Problem 20}\par
A taxi fare consists of a fixed fee of \(6\) units and \(2.5\) units for every kilometre travelled. Define a function \(C(d)\), state a meaningful domain, compute \(C(8)\), and explain whether the function is linear or affine.
\end{enumerate}

\section{Problem Set 9: Limits and Continuity}

\begin{enumerate}
\item \textbf{Problem 1}\par
Explain in words the difference between
\[
f(a)
\quad\text{and}\quad
\lim_{x\to a}f(x).
\]

\item \textbf{Problem 2}\par
Compute
\[
\lim_{n\to\infty}\frac{1}{n},
\qquad
\lim_{n\to\infty}\frac{5n+2}{2n-1}.
\]

\item \textbf{Problem 3}\par
Compute
\[
\lim_{n\to\infty}\frac{3n^2-n+1}{2-n^2}.
\]

\item \textbf{Problem 4}\par
Determine whether each sequence converges:
\[
a_n=(-1)^n,
\qquad
b_n=\frac{(-1)^n}{n},
\qquad
c_n=2+\frac{3}{n}.
\]

\item \textbf{Problem 5}\par
Compute
\[
\lim_{n\to\infty}\left(\sqrt{n^2+4n}-n\right)
\]
by rationalizing.

\item \textbf{Problem 6}\par
Compute by direct substitution:
\[
\lim_{x\to2}(3x^2-x+4).
\]

\item \textbf{Problem 7}\par
Compute
\[
\lim_{x\to3}\frac{x^2-9}{x-3}.
\]

\item \textbf{Problem 8}\par
Compute
\[
\lim_{x\to9}\frac{\sqrt{x}-3}{x-9}.
\]

\item \textbf{Problem 9}\par
Let
\[
f(x)=
\begin{cases}
2x+1,&x<2,\\
x^2-1,&x\ge2.
\end{cases}
\]
Find the left-hand limit, right-hand limit, two-sided limit, and value \(f(2)\).

\item \textbf{Problem 10}\par
Compute
\[
\lim_{x\to0^-}\frac{|x|}{x},
\qquad
\lim_{x\to0^+}\frac{|x|}{x},
\]
and decide whether the two-sided limit exists.

\item \textbf{Problem 11}\par
Compute
\[
\lim_{x\to\infty}\frac{4x^2-3x+1}{2x^2+5}.
\]

\item \textbf{Problem 12}\par
Find all horizontal asymptotes of
\[
f(x)=\frac{3x-1}{x+4}.
\]

\item \textbf{Problem 13}\par
Determine the one-sided infinite limits of
\[
f(x)=\frac{1}{x-2}
\]
as \(x\to2^-\) and \(x\to2^+\). State the vertical asymptote.

\item \textbf{Problem 14}\par
Check whether
\[
f(x)=\frac{x^2-4}{x-2},\qquad x\ne2,
\]
is continuous at \(x=2\). If not, give a value of \(f(2)\) that makes it continuous.

\item \textbf{Problem 15}\par
Find \(k\) so that
\[
f(x)=
\begin{cases}
kx+1,&x<3,\\
x^2-2,&x\ge3
\end{cases}
\]
is continuous at \(x=3\).

\item \textbf{Problem 16}\par
Classify the discontinuity at the indicated point:
\begin{enumerate}
\item \(\dfrac{x^2-1}{x-1}\) at \(x=1\);
\item \(\dfrac{1}{x+2}\) at \(x=-2\);
\item \(\dfrac{|x|}{x}\) at \(x=0\);
\item \(\sin(1/x)\) at \(x=0\).
\end{enumerate}

\item \textbf{Problem 17}\par
Compute
\[
\lim_{x\to0}\frac{\sin 5x}{2x}.
\]

\item \textbf{Problem 18}\par
Compute
\[
\lim_{x\to0}\frac{\tan 3x}{\sin 2x}.
\]

\item \textbf{Problem 19}\par
Compute
\[
\lim_{x\to0}(1+4x)^{1/x}.
\]

\item \textbf{Problem 20}\par
For each limit, identify the most appropriate first method and then compute:
\[
\lim_{x\to1}\frac{x^3-1}{x-1},
\qquad
\lim_{x\to\infty}\frac{2x+5}{x^2+1},
\qquad
\lim_{x\to0}\frac{\sin x}{x},
\qquad
\lim_{x\to4}\frac{x-4}{\sqrt{x}-2}.
\]
\end{enumerate}

\section{Problem Set 10: Derivatives and Function Analysis}

\begin{enumerate}
\item \textbf{Problem 1}\par
For
\[
f(x)=x^2+1,
\]
form and simplify the difference quotient
\[
\frac{f(a+h)-f(a)}{h}.
\]

\item \textbf{Problem 2}\par
Use the limit definition to find the derivative of
\[
f(x)=3x+2.
\]

\item \textbf{Problem 3}\par
Find the tangent line to
\[
f(x)=x^2-1
\]
at \(x=2\).

\item \textbf{Problem 4}\par
Differentiate, stating the relevant domain:
\[
f(x)=x^7,\qquad
g(x)=x^{-3},\qquad
h(x)=\sqrt{x}.
\]

\item \textbf{Problem 5}\par
Differentiate:
\[
f(x)=4e^x-3\ln x+2\sin x-\cos x.
\]

\item \textbf{Problem 6}\par
Differentiate using the product rule:
\[
f(x)=x^3e^x.
\]

\item \textbf{Problem 7}\par
Differentiate using the quotient rule:
\[
f(x)=\frac{x^2+1}{x+2}.
\]

\item \textbf{Problem 8}\par
Differentiate using the chain rule:
\[
f(x)=(2x-3)^6.
\]

\item \textbf{Problem 9}\par
Differentiate:
\[
f(x)=e^{\sin(x^2)}.
\]
Identify all layers of composition.

\item \textbf{Problem 10}\par
Find all critical points of
\[
f(x)=x^3-6x^2+9x.
\]

\item \textbf{Problem 11}\par
Study the monotonicity of
\[
f(x)=x^2-4x+1
\]
using a sign chart for \(f'\).

\item \textbf{Problem 12}\par
Find and classify the local extrema of
\[
f(x)=x^3-3x.
\]

\item \textbf{Problem 13}\par
Find the absolute maximum and minimum of
\[
f(x)=x^2-2x
\]
on \([-1,3]\).

\item \textbf{Problem 14}\par
Determine the intervals of convexity and concavity of
\[
f(x)=x^3-6x.
\]

\item \textbf{Problem 15}\par
Find all inflection points of
\[
f(x)=x^4-4x^3.
\]
Remember that \(f''(a)=0\) gives only a candidate.

\item \textbf{Problem 16}\par
Use the second derivative test to classify the critical points of
\[
f(x)=x^4-4x^2.
\]
For any inconclusive point, use another method.

\item \textbf{Problem 17}\par
Find all vertical and horizontal asymptotes of
\[
f(x)=\frac{2x+1}{x-3}.
\]

\item \textbf{Problem 18}\par
Perform a complete analysis of
\[
f(x)=x^3-3x^2.
\]
Include domain, zeros, behavior at infinity, monotonicity, extrema, convexity, and inflection points.

\item \textbf{Problem 19}\par
Perform a complete analysis of
\[
f(x)=\frac{x}{x-1}.
\]
Include domain, asymptotes, monotonicity, and a coherent description of the graph.

\item \textbf{Problem 20}\par
A quantity is modelled by
\[
Q(t)=t^3-6t^2+9t+4,\qquad t\ge0.
\]
Find \(Q'(t)\), determine when the quantity is increasing or decreasing, identify local extrema, and explain what the sign of \(Q'(t)\) means in a modelling context.
\end{enumerate}

\section{Problem Set 11: Integrals}

\begin{enumerate}
\item \textbf{Problem 1}\par
Find all antiderivatives of
\[
f(x)=6x^2-4.
\]

\item \textbf{Problem 2}\par
Compute, on appropriate intervals,
\[
\int x^5\,dx,
\qquad
\int x^{-4}\,dx,
\qquad
\int \sqrt{x}\,dx.
\]

\item \textbf{Problem 3}\par
Compute
\[
\int \left(3e^x-\frac{2}{x}+4\cos x\right)\,dx.
\]

\item \textbf{Problem 4}\par
Rewrite and integrate:
\[
\int \frac{\sqrt{x}}{x^3}\,dx.
\]

\item \textbf{Problem 5}\par
Use substitution to compute
\[
\int 2x(1+x^2)^5\,dx.
\]

\item \textbf{Problem 6}\par
Use substitution to compute
\[
\int \frac{e^x}{1+e^x}\,dx.
\]

\item \textbf{Problem 7}\par
Use integration by parts to compute
\[
\int xe^x\,dx.
\]

\item \textbf{Problem 8}\par
Compute
\[
\int x^2\cos x\,dx
\]
using integration by parts more than once.

\item \textbf{Problem 9}\par
Compute
\[
\int \frac{2x+3}{x^2+3x+5}\,dx.
\]

\item \textbf{Problem 10}\par
Compute
\[
\int \frac{x}{x^2-2x+5}\,dx
\]
by splitting the numerator into a derivative part and a remainder.

\item \textbf{Problem 11}\par
Compute
\[
\int \frac{x^2+2}{x-1}\,dx
\]
by polynomial division.

\item \textbf{Problem 12}\par
Compute
\[
\int \sin^4 x\cos x\,dx.
\]

\item \textbf{Problem 13}\par
Compute
\[
\int \cos^2 x\,dx
\]
using a trigonometric identity.

\item \textbf{Problem 14}\par
Evaluate
\[
\int_0^2 (3x+1)\,dx.
\]

\item \textbf{Problem 15}\par
Find the signed area represented by
\[
\int_{-2}^{2}x\,dx
\]
and compare it with the ordinary geometric area between \(y=x\) and the \(x\)-axis on \([-2,2]\).

\item \textbf{Problem 16}\par
Use properties of definite integrals to simplify:
\[
\int_0^5 f(x)\,dx-\int_2^5 f(x)\,dx,
\]
and
\[
\int_3^1 g(x)\,dx.
\]

\item \textbf{Problem 17}\par
Use the fundamental theorem of calculus to evaluate
\[
\int_1^4 \frac{1}{x}\,dx.
\]

\item \textbf{Problem 18}\par
Let
\[
A(x)=\int_0^x (t^2+1)\,dt.
\]
Find \(A'(x)\) and compute \(A(2)\).

\item \textbf{Problem 19}\par
Choose an appropriate integration method for each integral and justify the choice before computing:
\[
\int x\sin(x^2)\,dx,
\qquad
\int x\ln x\,dx,
\qquad
\int \frac{3x^2}{x^3+1}\,dx.
\]

\item \textbf{Problem 20}\par
Compute and verify by differentiation:
\[
\int \left(
\frac{2x}{1+x^2}
+x e^{x^2}
+\sin x\cos x
\right)\,dx.
\]
\end{enumerate}

\section{Problem Set 12: Applications of Calculus}

\begin{enumerate}
\item \textbf{Problem 1}\par
For each question, state whether the main tool is a derivative, a definite integral, optimisation, or a differential model:
\begin{enumerate}
\item find instantaneous velocity;
\item find total water discharged from a known flow rate;
\item choose dimensions giving the largest area;
\item model a population whose growth rate is proportional to its size.
\end{enumerate}

\item \textbf{Problem 2}\par
A rectangle has perimeter \(40\). Express its area as a function of one side and find the dimensions that maximise the area.

\item \textbf{Problem 3}\par
A rectangular enclosure is built beside a straight wall, so only three sides require fencing. If \(60\) metres of fencing are available, find the dimensions that maximise the area.

\item \textbf{Problem 4}\par
Squares of side \(x\) are cut from the corners of a \(20\text{ cm}\times30\text{ cm}\) sheet to form an open box. Build the volume function, determine the meaningful domain, and find the value of \(x\) that maximises the volume.

\item \textbf{Problem 5}\par
Find the area enclosed by
\[
y=4-x^2
\]
and the \(x\)-axis.

\item \textbf{Problem 6}\par
The curve
\[
y=\sqrt{x},\qquad 0\le x\le9,
\]
is rotated about the \(x\)-axis. Find the resulting volume.

\item \textbf{Problem 7}\par
Set up, but do not necessarily evaluate, the arc-length integral for
\[
y=x^2,\qquad 0\le x\le2.
\]

\item \textbf{Problem 8}\par
A particle has position
\[
s(t)=t^3-6t^2+9t,\qquad t\ge0.
\]
Find its velocity and acceleration.

\item \textbf{Problem 9}\par
For the particle in the previous problem, determine when it is at rest and on which intervals it is speeding up or slowing down.

\item \textbf{Problem 10}\par
A particle has acceleration
\[
a(t)=6t-4,
\]
initial velocity \(v(0)=3\), and initial position \(s(0)=2\). Find \(v(t)\) and \(s(t)\).

\item \textbf{Problem 11}\par
A spring obeys Hooke's law \(F(x)=40x\) newtons. Find the work required to stretch it from \(x=0.1\) m to \(x=0.25\) m.

\item \textbf{Problem 12}\par
Water flows into a tank at the non-negative rate
\[
r(t)=12-2t
\]
litres per minute for \(0\le t\le6\). How much water enters during the first four minutes?

\item \textbf{Problem 13}\par
The temperature over a six-hour interval is modelled by
\[
T(t)=18+2t-\frac{t^2}{3},\qquad 0\le t\le6.
\]
Find its average value on the interval.

\item \textbf{Problem 14}\par
A population satisfies
\[
N'(t)=0.04N(t),
\qquad
N(0)=500.
\]
Find \(N(t)\) and the doubling time.

\item \textbf{Problem 15}\par
A radioactive substance satisfies
\[
M'(t)=-0.08M(t),
\qquad
M(0)=120.
\]
Find \(M(t)\) and its half-life.

\item \textbf{Problem 16}\par
A quantity \(Q(t)\) has rate of change
\[
Q'(t)=3t^2-4t+1
\]
and \(Q(0)=7\). Find \(Q(t)\) and explain the role of the initial condition.

\item \textbf{Problem 17}\par
A closed cylindrical can must have fixed volume \(V\). Derive the relation between radius and height that minimises its surface area.

\item \textbf{Problem 18}\par
A factory is \(400\) m downstream on the opposite side of a river \(80\) m wide. Underwater cable costs three times as much per metre as land cable. Let \(x\) be the downstream distance from the point directly opposite the station to the landing point. Build the cost function and determine the optimal \(x\).

\item \textbf{Problem 19}\par
A car's velocity is
\[
v(t)=t^2-6t+8,\qquad 0\le t\le6.
\]
Find the signed displacement and the total distance travelled on the interval.

\item \textbf{Problem 20}\par
A company models daily profit by
\[
P(x)=-2x^3+30x^2-96x-50,
\qquad 0\le x\le10,
\]
where \(x\) is production in hundreds of units. Determine the production levels at which profit is locally maximal or minimal, compare all relevant endpoint and critical values, and interpret the economically best choice within the model.
\end{enumerate}
